Chain-of-thought reasoning is the dominant paradigm for improving large language model (LLM) performance, yet practitioners lack guidance on how reasoning trace length interacts with out-of-distribution (OOD) robustness.
We systematically investigate whether longer reasoning chains always improve generalization by varying token budgets across five levels in GPT-4.1 and evaluating on four benchmarks: MATH (in-distribution), GSM8K (easy OOD), MMLU-Pro (knowledge-heavy OOD), and MuSR (hard multi-step reasoning OOD).
We find that the length--accuracy relationship is \emph{task-dependent}: monotonically increasing for tasks within the model's capability range (MATH: 46\%$\rightarrow$78\%; MMLU-Pro: 63\%$\rightarrow$79\%), flat for easy tasks (GSM8K: ${\sim}$93\% across all budgets), and non-monotonic for tasks at the model's capability frontier (MuSR peaks at 22\% with moderate budgets before declining).
Critically, we discover that model-expressed confidence is well-calibrated in-distribution ($\rho{=}0.63$ on MATH) but completely uncalibrated on hard OOD tasks ($\rho{\approx}0$ on MuSR), undermining confidence-based adaptive length selection precisely where it is most needed.
An adaptive two-pass strategy achieves comparable accuracy to fixed long reasoning on MATH with 32\% fewer tokens, but fails catastrophically on MuSR due to systematic overconfidence.
These findings reveal that the ``more reasoning is better'' assumption breaks down at the frontier of model capability, with direct implications for efficient and robust LLM deployment.
